%-*-latex-*-
In order to transform one source string of text \verb!x[1..m]!
to a target string \verb!y[1..n]!,
we can perform various transformation operations.
Our goal is, given \verb!x! and \verb!y!,
to produce a series of transformations that change \verb!x! to \verb!y!.
We use an array \verb!z!
-- assumed to be large enough to hold all the characters it will need -- to hold
the intermediate results.
Initially, \verb!z! is empty, and at termination, we should have
\verb!z[j] = y[j]! for \verb!j = 1,2,...,n!.
We maintain current indices \verb!i! into \verb!x! and \verb!j! into \verb!z!,
and the operations are allowed to alter \verb!z! and these indices.
Initially, \verb!i = j = 1!.
We are required to examine every character in \verb!x! during the transformation, which
means that at the end of the sequence of transformation operations, we must have
\verb!i = m + 1!.

We may choose from among six transformation operations:

\begin{itemize}
  
\item
  Copy a character from \verb!x! to \verb!z! by setting \verb!z[j] = x[i]! and then incrementing both \verb!i!
  and \verb!j! . This operation examines \verb!x[i]!.

\item
  Replace a character from \verb!x! by another character \verb!c!, by setting \verb!z[j] = c!,
  and then
  incrementing both \verb!i! and \verb!j!.
  This operation examines \verb!x[i]!.

  \item
    Delete a character from \verb!x! by incrementing \verb!i! but leaving \verb!j! alone.
    This operation
    examines \verb!x[i]!.

  \item
    Insert the character \verb!c! into \verb!z! by setting \verb!z[j] = c! and then incrementing \verb!j!, but
    leaving \verb!i! alone.
    This operation examines no characters of \verb!x!.

  \item
    Twiddle (i.e., exchange) the next two characters by copying them from \verb!x! to \verb!z! but
    in the opposite order; we do so by setting \verb!z[j] = x[i + 1]! and \verb!z[j + 1] = x[i]!
    and then setting \verb!i = i + 2! and \verb!j = j + 2!. This operation examines \verb!x[i]! 
    and \verb!x[i + 1]!.

  \item
    Kill the remainder of \verb!x! by setting \verb!i = m + 1!.
    This operation examines all characters
    in \verb!x! that have not yet been examined.
    This operation, if performed, must
    be the final operation.
  \end{itemize}

  As an example, one way to transform the source string algorithm to the target
  string \verb!altruistic! is to use the following sequence of operations, where the
  underlined characters are \verb!x[i]! and \verb!z[j]! after the operation:


  \begin{Verbatim}[frame=single,commandchars=\\\{\}]
Operation         x            z
initial strings   \underline{a}lgorithm    _
copy              a\underline{l}gorithm    a_
copy              al\underline{g}orithm    al_
replace by t      alg\underline{o}rithm    alt_
delete            algo\underline{r}ithm    alt_
copy              algor\underline{i}thm    altr_
insert u          algor\underline{i}thm    altru_
insert i          algor\underline{i}thm    altrui_
insert s          algor\underline{i}thm    altruis_
twiddle           algorit\underline{h}m    altruisti_
insert c          algorit\underline{h}m    altruistic_
kill              algorithm_   altruistic_
\end{Verbatim}

Note that there are several other sequences of
transformation operations that transform
\verb!algorithm! to \verb!altruistic!.

Each of the transformation operations has an associated cost.
The cost of an operation depends on the specific application,
but we assume that each operation's
cost is a constant that is known to us.
We also assume that the individual costs of
the copy and replace operations are less than the combined costs of the delete and
insert operations; otherwise, the copy and replace operations would not be used.
The cost of a given sequence of transformation operations is the sum of the costs
of the individual operations in the sequence. For the sequence above, the cost of
transforming \verb!algorithm! to \verb!altruistic! is
\begin{align*}
  &3 \cdot \text{cost(copy)} + \text{cost(replace)} + \text{cost(delete)} + 4\cdot \text{cost(insert)} \\
  &\hspace{1cm} + \text{cost(twiddle)} + \text{cost(kill)}
\end{align*}

\begin{tightlist}
\item[(a)]
Given two sequences \verb!x[1..m]! and \verb!y[1..n]!
and set of transformation-operation
costs, the \defterm{edit distance} from x to y is the cost of the least expensive operation
sequence that transforms \verb!x! to \verb!y!. Describe a dynamic-programming algorithm
that finds the edit distance from \verb!x[1..m]! to \verb!y[1..n]! and prints an optimal operation
sequence. Analyze the running time and space requirements of your
algorithm.
Write a C++ program implementing your algorithm.
Here is a sample execution:
\begin{Verbatim}[frame=single,commandchars=\\\{\}]
\underline{algorithm}
\underline{altruistic}
\underline{1 2 3 1 2 5}
19
initial strings   >algorithm    
copy              a>lgorithm    a
copy              al>gorithm    al
replace by t      alg>orithm    alt
delete            algo>rithm    alt
copy              algor>ithm    altr
insert u          algor>ithm    altru
insert i          algor>ithm    altrui
insert s          algor>ithm    altruis
twiddle           algorit>hm    altruisti
insert c          algorit>hm    altruistic
kill              algorithm>    altruistic
\end{Verbatim}
The first two inputs are the strings \verb!x! and \verb!z!.
The next 6 inputs (integers) are the cost of copy, replace, delete, insert, twiddle, kill.
The first output is the cost.
The rest describe the operations.
\end{tightlist}

The edit-distance problem generalizes the problem of aligning two DNA sequences.
There are several
methods for measuring the similarity of two DNA sequences by aligning them.
One such method to align two sequences \verb!x! and \verb!y! consists of inserting spaces at
arbitrary locations in the two sequences (including at either end) so that the resulting
sequences \verb!x'! and \verb!y'! have the same length but do not have a space in the same
position (i.e., for no position \verb!j! are both \verb!x'[j]! and \verb!y'[j]! a space.) Then we assign a
\lq\lq score" to each position. Position \verb!j! receives a score as follows:
\begin{tightlist}
\li $+1$ if \verb!x'[j]! $=$ \verb!y'[j]! and neither is a space,
\li $-1$ if \verb!x'[j]! $\neq$ \verb!y'[j]! and neither is a space,
\li $-2$ if either \verb!x'[j]! or \verb!y'[j]! is a space.
\end{tightlist}
The score for the alignment is the sum of the scores of the individual positions. For
example, given the sequences \verb!x! $=$ \verb!GATCGGCAT! and \verb!y! $=$ \verb!CAATGTGAATC!, one
alignment is
\newpage
\begin{Verbatim}
    G ATCG GCAT
    CAAT GTGAATC
    -*++*+*+-++*
\end{Verbatim}
A \verb!+! under a position indicates a score of $+1$ for that position, a \verb!-! indicates a score
of $-1$, and a \verb!*! indicates a score of $-2$, so that this alignment has a total score of
$6 \cdot 1 - 2 \cdot 1 - 4 \cdot 2 = -4$.

\begin{tightlist}
\item[(b)]
Explain how to cast the problem of finding an optimal alignment as an edit
distance problem using a subset of the transformation operations copy, replace,
delete, insert, twiddle, and kill.
\end{tightlist}
